% =============================================================
%  Chapitre 6 — Tests et validation
% =============================================================
\chapter{Tests et validation}

\section*{Introduction}
La qualité d'un logiciel se mesure à la rigueur de sa stratégie de tests.
Ce chapitre présente la stratégie adoptée pour valider l'ensemble du système
Telnet Test Manager~: tests unitaires, tests d'intégration, tests
fonctionnels et résultats du pipeline CI/CD.

% -----------------------------------------------------------------
\section{Stratégie de tests}
% -----------------------------------------------------------------

\subsection{Pyramide de tests}

La stratégie de tests suit la \textbf{pyramide de tests} classique qui
préconise un grand nombre de tests unitaires rapides à la base, un nombre
modéré de tests d'intégration au milieu, et un nombre restreint de tests
fonctionnels (end-to-end) au sommet.

\begin{figure}[H]
  \centering
  \begin{tikzpicture}
    % Niveau 1 : Tests unitaires (base)
    \fill[bleuPFE!60] (0,0) -- (6,0) -- (5,1.5) -- (1,1.5) -- cycle;
    \node[white,font=\small\bfseries] at (3,0.75) {Tests Unitaires};
    \node[font=\scriptsize,white] at (3,0.35) {Jest — Use Cases, Workers, Middlewares};

    % Niveau 2 : Tests d'intégration
    \fill[bleuPFE!35] (1,1.5) -- (5,1.5) -- (4,3) -- (2,3) -- cycle;
    \node[white,font=\small\bfseries] at (3,2.25) {Tests d'intégration};
    \node[font=\scriptsize,white] at (3,1.85) {API REST + MongoDB};

    % Niveau 3 : Tests fonctionnels
    \fill[bleuPFE!20] (2,3) -- (4,3) -- (3.2,4.5) -- (2.8,4.5) -- cycle;
    \fill[bleuPFE!20] (2.8,4.5) -- (3.2,4.5) -- (3,5) -- cycle;
    \node[bleuPFE,font=\small\bfseries] at (3,3.75) {Tests E2E};
    \node[font=\scriptsize,bleuPFE!70] at (3,3.35) {Interface utilisateur};

    % Étiquettes de quantité
    \node[font=\scriptsize,color=bleuPFE] at (-1.5,0.75)
      {\textbf{Nombreux}};
    \node[font=\scriptsize,color=bleuPFE] at (-1.5,2.25)
      {\textbf{Modérés}};
    \node[font=\scriptsize,color=bleuPFE] at (-1.5,3.75)
      {\textbf{Peu}};

    % Axe vitesse
    \draw[->,>=stealth,gray!50] (-0.5,0) -- (-0.5,5.2)
      node[above,font=\scriptsize,gray] {Vitesse};

  \end{tikzpicture}
  \caption{Pyramide de tests adoptée pour le projet}
  \label{fig:test_pyramid}
\end{figure}

% -----------------------------------------------------------------
\section{Tests unitaires}
% -----------------------------------------------------------------

\subsection{Framework utilisé~: Jest}

Les tests unitaires sont rédigés avec le framework \textbf{Jest} (v29).
Chaque Use Case et chaque module critique dispose de sa suite de tests.
Les dépendances externes (MongoDB, réseau Telnet) sont mockées.

\subsection{Tests du module d'authentification}

\begin{lstlisting}[language=javascript,
  caption={Tests unitaires — LoginUseCase}]
describe('LoginUseCase', () => {

  const mockUserRepo = {
    findByUsername: jest.fn(),
  };
  const loginUseCase = new LoginUseCase(mockUserRepo);

  test('devrait retourner un JWT valide pour des identifiants corrects',
    async () => {
      const hashedPwd = await bcrypt.hash('motdepasse123', 10);
      mockUserRepo.findByUsername.mockResolvedValue({
        id: 1, username: 'testuser',
        password: hashedPwd, role: 'engineer',
        statut: 'actif',
      });

      const result = await loginUseCase.execute(
        'testuser', 'motdepasse123'
      );

      expect(result).toHaveProperty('token');
      expect(result.user.username).toBe('testuser');
    }
  );

  test('devrait rejeter des identifiants incorrects', async () => {
    mockUserRepo.findByUsername.mockResolvedValue({
      password: await bcrypt.hash('correct', 10), statut: 'actif',
    });

    await expect(
      loginUseCase.execute('user', 'mauvaismdp')
    ).rejects.toThrow('Identifiants invalides');
  });

  test('devrait rejeter un compte inactif', async () => {
    mockUserRepo.findByUsername.mockResolvedValue({
      password: await bcrypt.hash('mdp', 10), statut: 'inactif',
    });

    await expect(
      loginUseCase.execute('user', 'mdp')
    ).rejects.toThrow('Compte inactif');
  });
});
\end{lstlisting}

\subsection{Tests du Worker Thread Telnet}

\begin{lstlisting}[language=javascript,
  caption={Tests unitaires — Moteur Telnet (mock telnet-client)}]
jest.mock('telnet-client', () => {
  return jest.fn().mockImplementation(() => ({
    connect: jest.fn().mockResolvedValue(true),
    exec:    jest.fn().mockResolvedValue('PING 8.8.8.8 ... OK'),
    end:     jest.fn().mockResolvedValue(true),
  }));
});

describe('testWorker — mode single', () => {
  test('devrait exécuter une commande et valider la réponse',
    async () => {
      const result = await runSingleTest({
        slot:    { adresse: '192.168.5.1', port: 23 },
        command: {
          type: 'single',
          command: 'ping -c 1 8.8.8.8',
          expectedResponse: 'OK',
        },
      });

      expect(result.success).toBe(true);
      expect(result.response).toContain('OK');
    }
  );
});
\end{lstlisting}

% -----------------------------------------------------------------
\section{Tests d'intégration}
% -----------------------------------------------------------------

\subsection{Tests de l'API REST avec Supertest}

Les tests d'intégration utilisent \textbf{Supertest} pour simuler des
requêtes HTTP réelles vers l'API Express, avec une base MongoDB de test
(base isolée \texttt{telnet-test-db})~:

\begin{lstlisting}[language=javascript,
  caption={Tests d'intégration — Endpoint POST /login}]
describe('POST /login', () => {
  beforeAll(async () => {
    await connectDatabase(process.env.MONGO_URI_TEST);
    // Création d'un utilisateur de test
    await UserModel.create({
      id: 99, username: 'testuser',
      password: await bcrypt.hash('test1234', 10),
      role: 'engineer', statut: 'actif',
    });
  });

  afterAll(async () => {
    await UserModel.deleteMany({ username: 'testuser' });
    await mongoose.connection.close();
  });

  it('retourne 200 et un JWT pour des identifiants valides',
    async () => {
      const res = await request(app)
        .post('/login')
        .send({ username: 'testuser', password: 'test1234' });

      expect(res.status).toBe(200);
      expect(res.body).toHaveProperty('token');
      expect(res.body.user.role).toBe('engineer');
    }
  );

  it('retourne 401 pour un mot de passe incorrect', async () => {
    const res = await request(app)
      .post('/login')
      .send({ username: 'testuser', password: 'mauvais' });

    expect(res.status).toBe(401);
  });

  it('retourne 429 après 10 tentatives en 15 minutes',
    async () => {
      for (let i = 0; i < 10; i++) {
        await request(app).post('/login')
          .send({ username: 'x', password: 'x' });
      }
      const res = await request(app).post('/login')
        .send({ username: 'x', password: 'x' });
      expect(res.status).toBe(429);
    }
  );
});
\end{lstlisting}

% -----------------------------------------------------------------
\section{Tests fonctionnels (end-to-end)}
% -----------------------------------------------------------------

\subsection{Scénarios de tests manuels}

En l'absence d'un framework E2E (Playwright/Cypress) dans cette version,
les tests fonctionnels ont été réalisés manuellement selon des scénarios
documentés. Les principaux scénarios validés sont~:

\begin{table}[H]
  \centering
  \caption{Scénarios de tests fonctionnels — résultats}
  \label{tab:e2e_tests}
  \begin{tabularx}{\textwidth}{L{0.8cm} L{5cm} L{3cm} C{2cm}}
    \toprule
    \textbf{ID} & \textbf{Scénario} & \textbf{Précondition} & \textbf{Résultat} \\
    \midrule
    ST01 & Connexion avec identifiants valides &
           Utilisateur actif en BD & \textcolor{vertCode}{\textbf{PASS}} \\
    ST02 & Connexion avec mot de passe erroné &
           --- & \textcolor{vertCode}{\textbf{PASS}} \\
    ST03 & Accès à une route protégée sans JWT &
           --- & \textcolor{vertCode}{\textbf{PASS}} \\
    ST04 & Création d'un Poste (CRUD) &
           Rôle ingénieur & \textcolor{vertCode}{\textbf{PASS}} \\
    ST05 & Création d'un Slot avec adresse IP &
           Poste + Produit existants & \textcolor{vertCode}{\textbf{PASS}} \\
    ST06 & Exécution d'un test commande \texttt{single} &
           Slot accessible, équipement actif & \textcolor{vertCode}{\textbf{PASS}} \\
    ST07 & Exécution d'un test séquence (3 étapes) &
           Slot accessible & \textcolor{vertCode}{\textbf{PASS}} \\
    ST08 & Résultats en temps réel via WebSocket &
           Test en cours & \textcolor{vertCode}{\textbf{PASS}} \\
    ST09 & Arrêt d'un test en cours &
           Test PENDING & \textcolor{vertCode}{\textbf{PASS}} \\
    ST10 & Tests multi-slots (3 slots en parallèle) &
           3 slots configurés & \textcolor{vertCode}{\textbf{PASS}} \\
    ST11 & Génération d'un rapport &
           Tests historisés & \textcolor{vertCode}{\textbf{PASS}} \\
    ST12 & Consultation du journal d'audit &
           Rôle admin & \textcolor{vertCode}{\textbf{PASS}} \\
    ST13 & Changement de langue (FR → EN) &
           --- & \textcolor{vertCode}{\textbf{PASS}} \\
    ST14 & Accès technicien à une route admin &
           JWT technicien & \textcolor{vertCode}{\textbf{PASS}} \\
    \bottomrule
  \end{tabularx}
\end{table}

% -----------------------------------------------------------------
\section{Résultats du pipeline CI/CD}
% -----------------------------------------------------------------

\subsection{Exécution du pipeline}

Chaque poussée sur la branche \texttt{main} déclenche les trois jobs du
pipeline. Le tableau suivant résume les durées et statuts observés lors
des dernières exécutions~:

\begin{table}[H]
  \centering
  \caption{Résultats du pipeline GitHub Actions}
  \label{tab:cicd_results}
  \begin{tabularx}{\textwidth}{L{3cm} L{3.5cm} C{2.5cm} C{2.5cm}}
    \toprule
    \textbf{Job} & \textbf{Étapes clés} &
    \textbf{Durée} & \textbf{Statut} \\
    \midrule
    ci-backend & Checkout + npm ci + audit + tests &
      $\sim$1 min 45 s & \textcolor{vertCode}{\textbf{Success}} \\
    ci-frontend & Checkout + npm ci + audit + build &
      $\sim$2 min 30 s & \textcolor{vertCode}{\textbf{Success}} \\
    deploy & Build Docker + Helm upgrade + rollout &
      $\sim$6 min 20 s & \textcolor{vertCode}{\textbf{Success}} \\
    \bottomrule
    \multicolumn{2}{l}{\textbf{Durée totale}} &
      \multicolumn{2}{l}{$\sim$10 min 35 s} \\
    \bottomrule
  \end{tabularx}
\end{table}

\subsection{Vérification de l'état des pods après déploiement}

\begin{lstlisting}[language=yaml,
  caption={Sortie de \texttt{kubectl get pods -n telnet-app} après déploiement}]
NAME                                    READY   STATUS    RESTARTS   AGE
telnet-app-backend-7d4f9c8b6f-xk2p9    1/1     Running   0          4m12s
telnet-app-backend-7d4f9c8b6f-m8tn7    1/1     Running   0          3m55s
telnet-app-frontend-5c9d7b4f8d-qw3r1   1/1     Running   0          4m08s
telnet-app-mongodb-6b8f7c5d9b-nj4t8    1/1     Running   0          4m20s
\end{lstlisting}

% -----------------------------------------------------------------
\section{Tableau de synthèse de validation}
% -----------------------------------------------------------------

\begin{table}[H]
  \centering
  \caption{Synthèse de la validation par rapport aux besoins}
  \label{tab:synthese_validation}
  \begin{tabularx}{\textwidth}{C{0.8cm} L{5cm} L{3cm} C{2cm}}
    \toprule
    \textbf{BF} & \textbf{Besoin} & \textbf{Méthode de test} & \textbf{Statut} \\
    \midrule
    BF-01 & Authentification JWT & Tests unitaires + intégration &
      \textcolor{vertCode}{Validé} \\
    BF-05/06/07/08 & CRUD Configuration & Tests intégration &
      \textcolor{vertCode}{Validé} \\
    BF-09/10/11 & Gestion commandes Telnet & Tests intégration &
      \textcolor{vertCode}{Validé} \\
    BF-12 & Exécution test Telnet & Tests unitaires + fonctionnels &
      \textcolor{vertCode}{Validé} \\
    BF-13 & Modes single/séquence/monitoring & Tests fonctionnels &
      \textcolor{vertCode}{Validé} \\
    BF-14 & Tests multi-slots parallèles & Tests fonctionnels &
      \textcolor{vertCode}{Validé} \\
    BF-17/18/19 & WebSocket temps réel & Tests fonctionnels &
      \textcolor{vertCode}{Validé} \\
    BF-20 & Historique des tests & Tests intégration &
      \textcolor{vertCode}{Validé} \\
    BF-21/22 & Génération de rapports & Tests fonctionnels &
      \textcolor{vertCode}{Validé} \\
    BF-23/24 & Gestion utilisateurs RBAC & Tests unitaires + intégration &
      \textcolor{vertCode}{Validé} \\
    BF-25 & Journal d'audit & Tests fonctionnels &
      \textcolor{vertCode}{Validé} \\
    BF-26 & Statistiques admin & Tests fonctionnels &
      \textcolor{vertCode}{Validé} \\
    BF-27/28 & Internationalisation & Tests fonctionnels &
      \textcolor{vertCode}{Validé} \\
    --- & Pipeline CI/CD & Exécution GitHub Actions &
      \textcolor{vertCode}{Validé} \\
    --- & Déploiement Kubernetes & \texttt{kubectl rollout status} &
      \textcolor{vertCode}{Validé} \\
    --- & Supervision Prometheus/Grafana & Consultation dashboard &
      \textcolor{vertCode}{Validé} \\
    \bottomrule
  \end{tabularx}
\end{table}

\section*{Conclusion}

Ce chapitre a présenté la stratégie de validation du projet Telnet Test
Manager. Les tests unitaires (Jest), les tests d'intégration (Supertest),
les tests fonctionnels manuels et la validation du pipeline CI/CD
confirment que l'ensemble des 27 besoins fonctionnels et des exigences
non fonctionnelles sont satisfaits. Le taux de réussite de 100\,\% sur les
14 scénarios E2E testés témoigne de la qualité de la réalisation.
